\chapter[1994 --- Gilman et al.]{1994: Alfred G. Gilman, Martin Rodbell}
\label{ch:1994_Gilman}

\section*{Main Contribution}

\textit{Discovered G proteins and their role in signal transduction, explaining how cell surface receptors communicate with intracellular targets---a universal mechanism of cell communication.}

\section{Official Citation}

for their discovery of G proteins and the role of these proteins in signal transduction in cells

\section{Background and Historical Context}

Cells are constantly bombarded by signals from their environment---hormones, neurotransmitters, growth factors, and sensory stimuli. To respond appropriately to these signals, cells must possess elaborate molecular machinery to detect, amplify, and integrate incoming information. The discovery of G proteins and their role in signal transduction by Alfred G. Gilman and Martin Rodbell provided the conceptual and experimental framework for understanding this machinery. Their work, recognized with the 1994 Nobel Prize in Physiology or Medicine, revealed a universal mechanism by which cell surface receptors communicate with intracellular effector molecules, and its implications extend to virtually every aspect of physiology and pharmacology.

Alfred Goodman Gilman was born on July 1, 1941, in New Haven, Connecticut. He grew up in a family of physicians---his father, Louis S. Gilman, was a professor of pharmacology at Yale---and developed an early interest in both medicine and research. Gilman attended Trinity University in San Antonio, Texas, and then enrolled at Case Western Reserve University School of Medicine, where he earned his M.D.\ and Ph.D.\ in 1969. His doctoral work, under the supervision of Earl W. Sutherland (who had won the Nobel Prize in 1971 for his discovery of cyclic AMP), focused on the regulation of adenylyl cyclase by hormones. Gilman joined the faculty at the University of Virginia and later moved to the University of Texas Southwestern Medical Center in Dallas, where he conducted his Nobel Prize--winning research.

Martin Rodbell was born on December 1, 1925, in Baltimore, Maryland. He studied biochemistry at Johns Hopkins University, earning his B.A.\ in 1949 and his Ph.D.\ in 1954 under the supervision of Duane Bentley. After postdoctoral work with Claude ACC in Bethesda, Rodbell joined the National Institute of Arthritis and Metabolic Diseases (now the National Institute of Diabetes and Digestive and Kidney Diseases) at the National Institutes of Health, where he spent the remainder of his career. Rodbell's early research concerned the mechanism by which hormones stimulate the breakdown of glycogen in fat cells, and it was through these studies that he discovered the existence of G proteins.

The intellectual context of their work was shaped by the discovery of cyclic AMP (cAMP) as a second messenger by Earl Sutherland in the late 1950s. Sutherland showed that hormones such as epinephrine and glucagon stimulate the production of cAMP inside cells, and that cAMP in turn activates intracellular enzymes that mediate the hormone's metabolic effects. The enzyme responsible for cAMP production---adenylyl cyclase---was known to be located in the cell membrane and to be stimulated by hormones acting through cell surface receptors. However, the molecular mechanism by which receptors communicated with adenylyl cyclase was unknown. In the early 1970s, Rodbell proposed the existence of a ``transducer'' molecule---an intermediary protein that couples receptors to effector enzymes---based on his observation that the hormone-stimulated production of cAMP in fat cells required the presence of GTP (guanosine triphosphate).

\section{The Discovery/Contribution}

\subsection{Martin Rodbell and the GTP Requirement}

Rodbell's key insight came from a series of elegant experiments with fat cells (adipocytes) conducted in the late 1960s and early 1970s. He showed that the hormone epinephrine stimulated the production of cAMP in fat cell membrane preparations, but that this stimulation required the presence of GTP. Moreover, the non-hydrolyzable GTP analog GMP-P(NH)P was even more effective than GTP itself, suggesting that GTP hydrolysis played a regulatory role. Rodbell proposed that a GTP-binding protein---which he initially called ``N'' for nucleotide-binding regulatory protein---served as the intermediary between the hormone receptor and adenylyl cyclase.

This proposal was met with skepticism by many researchers, who questioned whether such a protein existed. However, Rodbell's experimental findings were robust and reproducible, and they pointed to a fundamental principle of signal transduction that would soon be confirmed and extended by Gilman.

\subsection{Alfred Gilman and the Identification of Gs}

Gilman took a biochemical approach to the problem, attempting to purify and characterize the GTP-binding protein that Rodbell had postulated. In the late 1970s, Gilman and his colleague Henry Bourne used a clever strategy: they identified cells that lacked a functional adenylyl cyclase response---a mutant cell line called CYC$^{-}$ (later renamed UNC)---and showed that these cells lacked a specific GTP-binding protein. By comparing normal cells with mutant cells, Gilman was able to identify the missing protein, which he named Gs (stimulatory G protein).

In 1980, Gilman's laboratory achieved the purification of Gs and demonstrated that it consisted of two subunits: an $\alpha$ subunit (G$\alpha$s) that binds GTP and activates adenylyl cyclase, and a $\beta\gamma$ subunit complex. GTP binding to the $\alpha$ subunit causes a conformational change that activates the protein and enables it to stimulate adenylyl cyclase. Hydrolysis of GTP to GDP by an intrinsic GTPase activity in the $\alpha$ subunit inactivates the protein, terminating the signal. This cycle of GTP binding, activation, hydrolysis, and inactivation constitutes the molecular switch that underlies G protein--mediated signaling.

\subsection{The G Protein Superfamily}

The discovery of Gs was followed by the identification of additional G proteins. Gilman, Rodbell, and other researchers discovered Gi (inhibitory G protein), which inhibits adenylyl cyclase; Gq, which activates phospholipase C; and transducin, which mediates phototransduction in the retina. These proteins share a common structural and functional motif: they are heterotrimeric (consisting of $\alpha$, $\beta$, and $\gamma$ subunits), they bind and hydrolyze GTP, and they serve as molecular switches that relay signals from cell surface receptors to intracellular effector molecules.

The human genome encodes 16 different G$\alpha$ subunits, 5 G$\beta$ subunits, and 12 G$\gamma$ subunits, providing an enormous combinatorial diversity of G protein signaling pathways. This diversity enables cells to respond to hundreds of different extracellular signals in a highly specific and finely tuned manner.

\subsection{G Protein-Coupled Receptors}

The discovery of G proteins also revealed the identity and mechanism of G protein-coupled receptors (GPCRs), a superfamily of cell surface receptors that represent the largest and most diverse group of membrane receptors in the human genome. The human genome encodes approximately 800 GPCRs, which mediate responses to hormones, neurotransmitters, chemokines, light, odorants, and many other stimuli. GPCRs are seven-transmembrane-domain proteins that, upon binding their ligand, undergo a conformational change that enables them to interact with and activate G proteins on the intracellular face of the membrane.

\section{Impact on Modern Medicine}

The impact of Gilman and Rodbell's discoveries on modern medicine is enormous. G protein--coupled receptors are the targets of approximately 34\% of all FDA-approved drugs, making them a major drug target class in pharmacology. Drugs acting on GPCRs include beta-blockers (for hypertension and heart failure), antihistamines (for allergies), opioids (for pain), antipsychotics (for schizophrenia), and many others.

\subsection{Pharmacological Implications}

Understanding the molecular mechanism of GPCR signaling has enabled the rational design of drugs that target specific receptor subtypes and signaling pathways. The development of selective agonists, antagonists, and inverse agonists has improved the efficacy and reduced the side effects of many medications. Biased agonists---drugs that selectively activate certain signaling pathways downstream of a receptor while sparing others---represent a new frontier in pharmacology that is directly informed by the G protein signaling framework.

\subsection{Disease Mechanisms}

Mutations in G proteins or GPCRs cause a range of diseases. Activating mutations in the G$\alpha$s gene cause McCune--Albright syndrome and certain forms of pituitary tumors. Inactivating mutations cause pseudohypoparathyroidism. Mutations in rhodopsin, a GPCR in the retina, cause retinitis pigmentosa. Understanding these disease mechanisms has guided the development of targeted therapies.

\subsection{Biased Signaling and New Drug Targets}

The recognition that GPCRs can signal through multiple pathways---G protein--dependent and G protein--independent (via $\beta$-arrestin)---has opened new avenues for drug discovery. The concept of biased agonism, in which drugs selectively activate beneficial pathways while avoiding harmful ones, is being applied to the development of next-generation opioids with reduced addiction potential and to new treatments for heart failure and metabolic diseases.

\section{Legacy}

Alfred Gilman became dean of the University of Texas Southwestern Medical Center and continued to study G protein signaling and its role in cancer. He received numerous honors, including the Albert Lasker Basic Medical Research Award in 1989. Gilman was diagnosed with bladder cancer and passed away on December 23, 2015, at the age of 74.

Martin Rodbell became scientific director of the National Institute of Environmental Health Sciences and continued to explore the broader implications of G protein signaling. He received the Albert Lasker Basic Medical Research Award in 1984 and was elected to the National Academy of Sciences. Rodbell died on December 7, 1998, at the age of 73.

Together, Gilman and Rodbell revealed the molecular logic of cell communication. Their discovery of G proteins and GPCR signaling has become one of a major frameworks in modern biology, underpinning our understanding of physiology, pharmacology, and disease. The ongoing exploration of GPCR signaling continues to yield new insights and therapeutic opportunities, ensuring that their legacy will endure for generations.


\section{Modern Developments}

Gilman and Rodbell's G protein work has led to modern pharmacology. Understanding of GPCR signaling has enabled development of drugs targeting G protein-coupled receptors, which are targets of ~34% of approved drugs. Understanding of GPCR structure has enabled structure-based drug design. Bias agonism---activating specific signaling pathways---has led to drugs with improved safety profiles. Understanding of GPCR desensitization has informed development of long-acting drugs.
